(1) 总电荷 \(Q=q+q-2q=0\)；以原点为参考，电偶极矩为 \(\mathbf p=q a\hat z+q(-a\hat z)-2q\cdot\mathbf0=\mathbf0\)。因单极与偶极系数都严格为零，远场电势的第一个非零多极项应从二阶四极开始。
(2) 对于 \(r>a\)，用 \(|\mathbf r\mp a\hat z|^{-1}=r^{-1}\sum_{\ell\ge0}(\pm a/r)^\ell P_\ell(\cos\theta)\)。两端等量正电荷相加消去奇阶，中心负电荷抵消零阶，得到 \(\Phi=2qa^2P_2(\cos\theta)/(4\pi\varepsilon_0r^3)+O(r^{-5})=qa^2(3\cos^2\theta-1)/(4\pi\varepsilon_0r^3)+O(r^{-5})\)。
(3) 在 \(\theta=0\) 轴向有 \(3\cos^2\theta-1=2>0\)，电势主项为正；赤道 \(\theta=\pi/2\) 有因子 \(-1\)，主项为负。径向电场 \(E_r=-\partial_r\Phi\)，因此主导按 \(r^{-4}\) 衰减，比偶极远场电场的 \(r^{-3}\) 快一个幂次。
(4) 保持电荷量固定而让 \(a\to sa\)，四极远场系数 \(qa^2\) 增至 \(s^2qa^2\)；下一项 \(\ell=4\) 的大小随 \(a^4/r^5\)，仍需 \(r\gg sa\) 才是小修正。这个缩放规律验证四极场确由电荷分离的二次长度矩控制。
